% =============================================================================
% CHAPTER 1 - INTRODUCTION
% =============================================================================

\chapter{INTRODUCTION}

\ifpdf
  \graphicspath{{introduction/figures/PNG}{introduction/figures/PDF}{introduction/figures}}
\else
  \graphicspath{{introduction/figures/EPS}{introduction/figures}}
\fi

\section{Background}

Over the past decade, large language models (LLMs) have grown from narrow text predictors into systems capable of multi-step reasoning, nuanced summarisation, and broad question answering~\citep{brown2020gpt3,openai2024gpt4technicalreport}. Models such as GPT-4 and the Llama family acquire these capabilities through exposure to vast text corpora, internalising linguistic patterns and factual associations as statistical regularities in their parameters~\citep{touvron2023llama,dubey2024llama3}. The same process that makes them broadly capable, however, makes them structurally unreliable on knowledge-intensive tasks: because knowledge is stored as implicit weights rather than as verifiable entries, a model can produce fluent, confident answers that are factually wrong — a problem known as hallucination~\citep{ji2023survey}. \citet{ji2023survey} document rates ranging from 15\% to over 60\% across summarisation, dialogue, and question-answering tasks. Scaling model size does not resolve this; larger models hallucinate differently, not less.

The root cause is architectural. At each generation step, the model selects the next token from a probability distribution conditioned only on its learned parameters and the tokens generated so far. There is no internal checkpoint against an external source of truth. If an incorrect token enters the sequence at step $t$, every subsequent step conditions on it, propagating the error coherently and fluently through the rest of the output~\citep{huang2023survey,ji2023survey}. This mechanism makes knowledge graph question answering (KGQA) a particularly demanding setting. A knowledge graph encodes facts as verified triplets — for example, (\textit{Marie Curie}, \textit{award\_received}, \textit{Nobel Prize in Physics}) — and answering complex questions requires traversing chains of these triplets in the correct order~\citep{bollacker2008freebase,lan2022survey}. On benchmarks such as WebQSP~\citep{yih2016websqp} and ComplexWebQuestions~\citep{talmor2018cwq}, a three-hop question can correspond to tens of thousands of structurally valid graph paths, and an unconstrained LLM must navigate this space using parametric memory alone, which is precisely the mechanism that leads to fabricated answers.

Retrieval-augmented generation (RAG) tries to address this by supplying relevant passages or subgraphs as context before generation~\citep{lewis2020rag,izacard2021leveraging}. This helps in practice, but not in principle. Providing context does not prevent the model from generating beyond it; the distribution over the full vocabulary is still computed autoregressively at every step, meaning any token retains a nonzero probability. \citet{luo2025-graph-constrained-reasoning} find that even when a knowledge graph is directly available, retrieval-augmented approaches still hallucinate reasoning steps in roughly one-third of cases. Soft context, however relevant, cannot guarantee faithfulness.

A structurally different remedy is constrained decoding. Rather than changing what the model sees, constrained decoding changes what it is allowed to generate. A constraint oracle computes a binary mask at each step, setting the log-probability of invalid tokens to $-\infty$ before the softmax so that their selection probability is exactly zero~\citep{geng2023grammarconstrained,willard2023outlines}. Graph-Constrained Reasoning (GCR)~\citep{luo2025-graph-constrained-reasoning} applies this principle to KGQA by pre-building a KG-Trie — a prefix tree encoding all valid graph paths within $L$ hops of the question entities — and using it as the oracle during beam-search decoding. GCR achieves 92.6 Hits@1 on WebQSP and verified 100\% structural faithfulness, meaning every generated token corresponds to a real fact in the graph.

A meaningful gap remains, however. GCR's trie is constructed once before generation, from graph topology and question entities only, without any consideration of what the question actually means or how far reasoning has progressed. The valid token set is therefore identical at the first generation step and the tenth, regardless of context. On multi-hop questions this makes the oracle excessively permissive: thousands of structurally valid but semantically irrelevant paths remain in the constraint set at every step, and the model must filter among them using parametric knowledge — reintroducing the internal-memory reliance that constrained decoding was designed to eliminate. Decoding on Graphs (DoG)~\citep{li2024-decoding-on-graphs} updates the trie step-by-step as generation proceeds, but expansion is guided by graph topology alone; question semantics play no role in determining which paths remain valid. Neither method narrows the constraint set in response to what the question is asking.

\section{Problem Statement}

Existing KG-constrained decoding frameworks, including GCR~\citep{luo2025-graph-constrained-reasoning} and DoG~\citep{li2024-decoding-on-graphs}, define valid tokens using only graph structure and question entities. Formally, the valid-token set at each step is $W_{\mathrm{val}} = f(G, E_q)$, which does not depend on the input question $q$ or the partial generation $y_{<t}$. Structural faithfulness is preserved, but contextual relevance is not: the oracle admits every path reachable in the graph, whether or not it points in the direction the question requires.

This creates two compounding problems. First, the constraint set is too broad — at deep hop depths, valid but irrelevant paths can number in the thousands~\citep{he2021improving,talmor2018cwq}. Second, the constraint set is static — it does not narrow as the model commits to a reasoning direction. Together, these properties force the LLM to do exactly what constrained decoding was meant to prevent: selecting among candidates on the basis of parametric knowledge, with all the hallucination risk that entails.

This project addresses that gap. We propose DCA-Trie, which extends the validity function to $W_{\mathrm{val}}(t) = f(G, E_q, q, y_{<t})$ — conditioning the valid set jointly on the graph, the question, and the partial generation at each step. The framework must satisfy three requirements: (i) preserve 100\% structural faithfulness, so every accepted token matches a verified KG path; (ii) reduce semantic irrelevance, measured by a new metric called the Semantic Irrelevance Ratio (SIR; Chapter~\ref{Chapter3}); and (iii) maintain high answer recall, with false negatives on gold answer paths below 5\% on the WebQSP validation set. We also note that RouterKGQA~\citep{yuan2026routerkgqa}, a concurrent method, applies semantic filtering at the answer-selection stage after generation. DCA-Trie differs by filtering at the token-validity stage during decoding, preventing semantically irrelevant paths from being considered at all rather than discarding them afterwards.

\section{Objectives of Thesis}

The aim of this project is to design, implement, and evaluate a Dynamic Context-Aware Trie (DCA-Trie) framework for knowledge graph-constrained large language model generation, in which the valid token constraint at each decoding step is conditioned jointly on the knowledge graph structure, input question semantics, and partially generated reasoning chain, improving answer accuracy on multi-hop KGQA benchmarks while maintaining 100\% structural faithfulness to the underlying knowledge graph. Specific objectives are as follows:

\begin{enumerate}[label=\roman*.]
  \item to characterise the permissiveness problem in existing KG-constrained decoding frameworks by defining and measuring the Semantic Irrelevance Ratio (SIR) of GCR's static KG-Tries on the WebQSP and CWQ benchmarks, stratified by hop depth;
  \item to design a semantic relevance scoring mechanism using pretrained sentence transformer embeddings~\citep{reimers2019sbert} that filters candidate KG paths by joint relevance to the input question and partial generation at inference time, without task-specific fine-tuning;
  \item to implement a static semantic filtering variant (DCA-Trie v1) that applies the relevance scorer at KG-Trie construction time, reducing trie size before generation begins while maintaining a false negative rate below 5\% on gold answer paths on the WebQSP validation set;
  \item to implement a step-wise dynamic expansion variant (DCA-Trie v2) in which the trie is updated after each committed triplet using the relevance scorer conditioned jointly on the question and partial generation $y_{<t}$, directly operationalising the $W_{\mathrm{val}}(t) = f(G, E_q, q, y_{<t})$ formulation;
  \item to evaluate DCA-Trie v1 and v2 against the GCR baseline and an unconstrained chain-of-thought (CoT) baseline on WebQSP and CWQ, measuring Hits@1, F1, structural faithfulness, SIR, and average trie size, with results stratified by question hop depth.
\end{enumerate}

\section{Tools and Facilities Used}

\subsection{Facilities}

The following resources were used in carrying out this project:

\begin{enumerate}[label=\roman*.]
  \item Google Colab Pro (A100, 40 GB VRAM) for all model inference and constrained decoding experiments.
  \item University library and online academic databases (ACM Digital Library, arXiv, Google Scholar, ACL Anthology) for literature review and citation verification.
  \item University internet infrastructure for cloud connectivity, dataset access, and retrieval of model checkpoints from the HuggingFace model hub.
  \item A shared team repository on GitHub for version control, experiment logging, and reproducibility of all reported results.
\end{enumerate}

\subsection{Software Tools}

Table~\ref{tab:implemented-software-tools} presents the software tools used in this project, together with their role in the DCA-Trie framework and the rationale for their selection.

\begin{table}[H]
  \centering
  \small
  \setlength{\tabcolsep}{4pt}
  \caption[\normalfont Implemented Software Tools]{Implemented Software Tools.}
  \label{tab:implemented-software-tools}
  \resizebox{\textwidth}{!}{%
    \begin{tabular}{|>{\raggedright\arraybackslash}p{0.20\textwidth}|>{\raggedright\arraybackslash}p{0.30\textwidth}|>{\raggedright\arraybackslash}p{0.35\textwidth}|}
      \hline
      Tool                                     & Use in Project                                                                                                    & Rationale                                                                                                                       \\
      \hline
      Python / HuggingFace Transformers        & Model inference pipeline; loading and running GCR's Llama-3.1-8B checkpoint; integrating the relevance scorer     & Provides reliable model loading, tokenisation, and generation utilities, enabling direct use of the released GCR checkpoint. \\
      \hline
      GCR Framework                            & Primary baseline; KG-Trie construction, graph-constrained decoding, and result reproduction on WebQSP and CWQ     & Provides a strong, publicly reproducible baseline and the trie-based decoding mechanism that DCA-Trie extends.                    \\
      \hline
      Sentence-transformers / all-MiniLM-L6-v2 & Semantic relevance scoring; computing cosine similarity between question-context and candidate KG path embeddings & Delivers fast sentence-level embeddings without fine-tuning, supporting the latency requirements of step-wise decoding.                          \\
      \hline
      PyTorch                                  & Deep learning framework for model inference and custom decoding logic                                & Supports implementation of custom decoding routines and integrates smoothly with HuggingFace.                          \\
      \hline
      Freebase / WebQSP and CWQ Datasets       & Knowledge graph source and KGQA benchmarks for evaluation                                                         & Standard resources in this research area, enabling direct comparison with prior work on identical evaluation splits.                                \\
      \hline
      Google Colab Pro (A100, 40 GB)           & GPU environment for all inference experiments; Llama-3.1-8B constrained decoding with beam search                 & Provides sufficient GPU memory for stable experiments without requiring local high-end hardware.                             \\
      \hline
      Visual Studio Code + Git / GitHub        & Development environment and version control                                                                       & Supports day-to-day development, team collaboration, and experiment reproducibility.                                           \\
      \hline
    \end{tabular}
  }
\end{table}

\section{Scope of Work}

This project designs and evaluates the DCA-Trie framework for KG-constrained LLM decoding on multi-hop KGQA over Freebase-grounded benchmarks (WebQSP and CWQ). All experiments use GCR's publicly released Llama-3.1-8B checkpoint, which includes their fine-tuned KG-specialised parameters; no additional model training or fine-tuning is performed. The semantic relevance scorer uses the all-MiniLM-L6-v2 sentence transformer without domain-specific adaptation. Our contribution is strictly the constraint oracle architecture and the semantic filtering mechanism.

Several extensions are explicitly out of scope. Generalisation to knowledge graphs other than Freebase, and application to tasks beyond KGQA (entity linking, structured information extraction), are identified as directions for future work and are not evaluated here. Formal proofs of faithfulness under dynamic pruning are beyond scope; we provide empirical measurement of the gold-answer false negative rate as evidence. The confidence-conditioned constraint tightening mechanism described in Chapter~\ref{Chapter3} is a theoretical extension that is not implemented. Production deployment, latency optimisation for online serving, and integration with proprietary knowledge graphs are engineering concerns outside the research scope of this project.

The primary contributions of this work are the DCA-Trie framework and the Semantic Irrelevance Ratio metric. Secondary contributions are the empirical characterisation of constraint permissiveness in existing frameworks and the ablation analysis isolating the contribution of semantic filtering from dynamic expansion. All code, experimental configurations, and curated evaluation data will be shared to support reproducibility.

\section{Organization of Work}

The remainder of this dissertation is structured as follows. Chapter~2 reviews the literature on LLM hallucination, multi-hop KGQA, constrained decoding, and retrieval-augmented generation, identifying the specific technical gaps that motivate DCA-Trie. Chapter~3 formalises the problem, defines the SIR metric, and details the design and implementation of both DCA-Trie variants within the GCR pipeline~\citep{luo2025-graph-constrained-reasoning}. Chapter~4 presents quantitative results on WebQSP and CWQ, including faithfulness measurements, SIR analysis, and ablation studies. Chapter~5 summarises the main contributions, discusses practical implications, and proposes directions for future work.
